\chapter{Design of a Frequent Batch Auction Engine}
\label{ch:design}

This chapter answers RQ1. It decomposes the engine into five design
blocks, enumerates the options available within each block, and derives
from the theory of \Cref{ch:foundations} the strategic behaviour each
option should be expected to induce. \Cref{sec:tradeoff} consolidates
the result into a trade-off matrix and \Cref{sec:casestudies} tests the
taxonomy against systems that have actually been deployed.

\section{The Engine as a Set of Design Blocks}
\label{sec:blocks}

A matching engine is a deterministic function from a stream of
timestamped messages to a stream of trades and book states. A continuous
engine applies that function incrementally: every message is processed
to completion before the next is read. A batch engine interposes an
accumulation phase, so the function is applied once per interval to a
set of messages rather than once per message.

That single change forces five decisions that a continuous engine never
has to make, because in a continuous engine each of them is implied by
serial processing.

\begin{enumerate}
  \item \textbf{The batch interval $\tau$.} How long the accumulation
  phase lasts and whether its end is predictable.
  \item \textbf{Information disclosure.} What participants observe
  during accumulation.
  \item \textbf{Order flow segregation.} Whether all orders compete in
  one pool.
  \item \textbf{Matching, clearing and rationing.} How the clearing
  price is computed and how the long side is allocated.
  \item \textbf{Settlement and residual handling.} How the cleared batch
  is settled and what happens to what did not clear.
\end{enumerate}

\Cref{fig:pipeline} shows how the blocks compose into a single batch
cycle. The remainder of the chapter treats each block in turn.

\begin{figure}[htbp]
\centering
\fbox{\begin{minipage}{0.92\textwidth}
\vspace{0.4em}
\centering
\ttfamily\small
\begin{tabular}{c}
\textbf{[1] Accumulation phase, duration $\tau$}\\
orders, cancels and modifies buffered; nothing executes\\
$\downarrow$ \quad \emph{governed by Block 1 (interval) and Block 2 (disclosure)}\\[0.5em]
\textbf{[2] Gate closes}\\
buffer frozen; late messages routed to batch $n+1$\\
$\downarrow$\\[0.5em]
\textbf{[3] Partition into auctions}\\
one pool, or maker/taker and retail/institutional partitions\\
$\downarrow$ \quad \emph{governed by Block 3 (segregation)}\\[0.5em]
\textbf{[4] Clearing}\\
solve for $\clr$ maximising executed volume\\
$\downarrow$ \quad \emph{governed by Block 4 (clearing)}\\[0.5em]
\textbf{[5] Rationing}\\
allocate the long side among eligible orders\\
$\downarrow$ \quad \emph{governed by Block 4 (rationing)}\\[0.5em]
\textbf{[6] Settlement and residual handling}\\
transfers executed; unmatched volume rolled over or cancelled\\
$\downarrow$ \quad \emph{governed by Block 5}\\[0.5em]
\textbf{[7] Publication} $\rightarrow$ next batch begins
\end{tabular}
\vspace{0.4em}
\end{minipage}}
\caption[The batch cycle and its design blocks]{One batch cycle of a
frequent batch auction engine, annotated with the design block that
governs each stage.}
\label{fig:pipeline}
\end{figure}
\section{Block 1: The Batch Interval}
\label{sec:interval}

\subsection{The Length of \texorpdfstring{$\tau$}{tau}}

The interval length is the parameter that determines how much of the
speed advantage the mechanism neutralises. Its effect is
straightforward: a speed differential of $\delta$ between two
participants matters only if it changes which batch their orders land
in, which happens only if their submission decisions fall within
$\delta$ of a boundary. For a uniformly distributed submission time the
probability of this is $\delta/\tau$. With $\delta$ on the order of
microseconds, which is the scale at which the races documented by
\citet{aquilina2022quantifying} are decided, and $\tau$ on the order of
100~milliseconds, the probability is on the order of $10^{-4}$ to
$10^{-5}$. Speed has not been prohibited; it has been made almost
worthless.

The cost is delay. Every participant waits, in expectation, $\tau/2$
before its order can execute, and $\tau$ in the worst case. For an
investor with a horizon of minutes this is negligible. For a market
maker managing inventory against a moving reference price it is a real
exposure, because the price can move during the interval while the
submitted order cannot be withdrawn in time to avoid the resulting
clearing.

The choice of the batch interval $\tau$ involves a trade-off: it must be
long enough to neutralise speed differences, but short enough to keep
prices current. Its appropriate value also depends on the liquidity of
the instrument.

\begin{itemize}
  \item \textbf{Lower bound:} $\tau$ must be comfortably larger than the
  differences in network latency across participants. Otherwise, faster
  participants will be more likely to reach the current batch in time,
  and the speed race will persist, only in a weaker form.
  \item \textbf{Upper bound:} $\tau$ must be short enough that new
  information is quickly reflected in prices. \citet{zoican2021speed}
  show that auction frequency does not improve market quality
  monotonically: very long intervals allow orders to become stale and
  raise the risk that the auction clears at a price the market has
  already moved away from.
  \item \textbf{Dependence on liquidity:} the right $\tau$ also depends
  on how much trading interest arrives in each batch. For a thinly
  traded instrument, a short interval may collect no matching buy and
  sell orders, leaving batches empty; the mechanism then behaves like a
  continuous market with a delay. The interval length should therefore
  be set together with the liquidity of the instrument, not
  independently of it.
\end{itemize}

\subsection{Fixed versus Randomised Closing}

If the closing time is known in advance, it becomes a coordination
point. A fast participant can hold its decision until the last feasible
moment, observe the maximum amount of information, and then submit or
cancel. This is boundary sniping, and it is the interval-scale
reappearance of the continuous-time race, with one difference: the
prize is not a single stale quote but participation in a batch whose
composition is now partly known.

\begin{itemize}
  \item \textbf{Fixed deterministic boundary.} Simple, auditable,
  trivially schedulable by participants, and compatible with
  wall-clock-aligned reporting. Maximally exposed to boundary sniping.
  \item \textbf{Randomised closing time.} The gate closes at a random
  time drawn from a distribution announced in advance, for example
  uniform on $[\tau_{\min}, \tau_{\max}]$. Waiting until the last
  moment now risks missing the batch entirely, which converts a
  dominant strategy into a gamble. The cost is that participants can no
  longer plan submission precisely, and that the randomisation source
  must be verifiable, which is a non-trivial requirement in a
  decentralised setting.
  \item \textbf{Volume- or event-triggered closing.} The auction closes
  when accumulated volume crosses a threshold, or on a market event.
  This adapts the interval to activity, which addresses the thin
  instrument problem, but makes the closing time endogenous to
  participant behaviour and therefore manipulable: a participant can
  trigger the auction by submitting volume.
  \item \textbf{Randomised extension.} A deterministic nominal close
  followed by a short random extension if late activity is detected.
  This is the design used in several equity closing auctions and is a
  compromise: predictable in the normal case, unpredictable exactly when
  predictability would be exploited.
\end{itemize}



\section{Block 2: Information Disclosure}
\label{sec:disclosure}

The disclosure regime determines what participants can condition on when
they submit. \Cref{tab:disclosure} sets out the four regimes introduced
in \Cref{sec:transparency} together with the behaviour each induces.

\begin{table}[htbp]
\centering
\caption[Disclosure regimes and induced behaviour]{Disclosure regimes
during the accumulation phase and their predicted behavioural
consequences.}
\label{tab:disclosure}
\small
\begin{tabularx}{\textwidth}{@{}p{2.9cm} L L@{}}
\toprule
\textbf{Regime} & \textbf{What is published during accumulation} &
\textbf{Predicted behaviour}\\
\midrule
Fully open book & All accumulated orders, continuously updated &
Strong within-batch price discovery; early submitters expose intent and
are exploited by later ones; incentive to submit as late as possible,
which compounds boundary sniping\\
Indicative clearing price & The price at which the batch would clear if
it closed now & Coordination on a common price estimate without exposing
individual orders; some intent leakage through price movement; the
regime of most equity opening and closing auctions\\
Aggregate imbalance & Net excess demand or supply, no price schedule &
Signals direction of pressure and attracts offsetting liquidity without
revealing individual limits; weaker price anchoring\\
Sealed bid & Nothing until the auction clears & No exploitation of
submitted orders and no incentive to delay for informational reasons;
clearing price uncertainty induces bid shading, which can reduce cleared
volume\\
\bottomrule
\end{tabularx}
\end{table}

Two points deserve emphasis because they are easy to miss.

First, disclosure and interval design are not independent. An open book
during accumulation strengthens the incentive to submit late, because
information accrues over the interval and a late submitter conditions on
more of it. An open book therefore raises the value of a randomised
close, whereas a sealed-bid design removes most of the informational
motive for waiting and makes a fixed boundary more defensible.

Second, opacity does not eliminate information leakage; it relocates it.
Under sealed bidding, the clearing price and volume of batch $n$ are
published before batch $n+1$ closes, so participants learn about
aggregate interest with a one-batch lag. As $\tau$ shrinks, the sequence
of published clearing prices approaches a continuous price signal, and
the difference between a sealed-bid FBA and a transparent one narrows in
informational terms even though it remains sharp within any single
batch.

\section{Block 3: Order Flow Segregation}
\label{sec:segregation}

\subsection{The Unified Pool}

In the default design, all orders enter one auction. This maximises the
set of potential counterparties and is the simplest to reason about and
to regulate. Its cost follows directly from \citet{glosten1985bid}: a
liquidity provider that cannot distinguish informed from uninformed flow
must quote a price reflecting the average information content of the
pool, so uninformed participants cross-subsidise the losses that
informed participants inflict.

\subsection{Segregation by Participant Class}

Splitting flow by participant class, typically retail against
institutional, allows a liquidity provider to quote a tighter price
against the subset it believes to be uninformed. This is the economic
logic behind retail wholesaling arrangements in equity markets, and it
transfers cleanly to a batch setting: run one auction per segment, or
one auction with segment-conditional eligibility.

The design questions it raises are mostly about classification. The
segmentation is only as good as the label, labels can be gamed by
routing institutional flow through retail channels, and segregation
raises fairness and best-execution questions because participants in
different segments receive different prices for the same asset at the
same instant. 

\subsection{Intent-Based Segregation: The Dual Flow Batch Auction}

Jump Trading's Dual Flow Batch Auction (DFBA) introduced by \citep{jump2025dfba} proposes partitions by
\emph{intent} rather than by participant identity. Every order is
labelled as maker flow, meaning liquidity supply, or taker flow, meaning
executable demand. Within each batch the engine then runs two
independent auctions:

\begin{itemize}
  \item maker buys against taker sells, and
  \item maker sells against taker buys,
\end{itemize}

both settled at a single uniform clearing price. Two structural
consequences follow. Market makers never trade against each other,
which removes the inter-maker crossing that in a unified pool produces
volume without providing liquidity to anyone who wanted it. And a fast
participant cannot use taker flow to pick off maker flow inside the
batch, because the batch clears simultaneously at one price and there is
no queue to win.

The design carries its own costs. Splitting the pool splits the
liquidity, so each of the two auctions clears against a thinner book
than a unified auction would. Intent labelling must be enforced rather
than declared, since a participant that can label itself a maker in one
batch and a taker in the next captures the benefits of both. And the
partition presumes that maker and taker are meaningful categories, which
in a mechanism that clears everything simultaneously is a stipulation
imposed by the engine rather than a fact about arrival timing.

\subsection{The Maker--Taker Distinction and Fee Design}

This last point generalises into a problem the FBA literature has not
resolved. In a continuous book, maker and taker are defined by arrival
timing relative to book state: the resting order is the maker, the
crossing order is the taker. That definition is unavailable in a
mechanism where all matched orders arrive during the same interval and
execute at the same instant at the same price. Yet fee schedules,
rebates, and a substantial part of exchange revenue depend on the
distinction.

Four ways to redefine it are available, none of them fully
satisfactory:

\begin{enumerate}
  \item \textbf{By declared intent}, as in DFBA. Operationally simple,
  but self-declared unless enforced by a separate eligibility regime.
  \item \textbf{By limit price aggressiveness.} Orders whose limit is
  far from the clearing price on the passive side are makers; orders
  that would have crossed a wide range of prices are takers. This is
  observable and hard to game without changing economic exposure, but it
  taxes participants for demanding certainty of execution.
  \item \textbf{By marginality.} The orders that set the clearing price
  are takers; inframarginal orders are makers. Theoretically appealing,
  because it is exactly the marginal order that determines the price,
  but it makes the fee a participant pays depend on the actions of
  everyone else in the batch, which is unpredictable ex ante.
  \item \textbf{By persistence across batches.} An order that has rested
  through several batches is a maker; a new order is a taker. This
  rewards standing liquidity, which is the economic function the rebate
  is supposed to compensate, and it is the only one of the four that
  measures something continuous rather than instantaneous.
\end{enumerate}

The choice is not merely accounting. A rebate schedule is an incentive,
and whichever definition the engine adopts becomes the behaviour it
subsidises.



\section{Block 4: Matching, Clearing and Rationing}
\label{sec:clearing}

This section turns the clearing problem of \Cref{sec:priceimprovement}
into an engine specification. It answers three questions left open
there: how the tie-breaking rule $\sigma$ works in practice, why every
order in a batch trades at the same price, and who is filled when the
two sides do not match exactly.

\subsection{Implementing the Clearing Rule}

In practice, prices can only take values on the tick grid, so the set
of volume-maximising prices $\mathcal{P}^{*}$ is a short run of
consecutive ticks. The engine finds it in a single pass over the
submitted limit prices, and then picks one of them:
\begin{equation}
  \clr = \sigma\bigl(\mathcal{P}^{*}\bigr).
  \label{eq:clearing}
\end{equation}
Exchanges implement $\sigma$ as a fixed sequence of tie-breakers:

\begin{enumerate}
  \item choose the price that leaves the smallest unmatched quantity;
  \item if still tied, choose the price on the side of the remaining
  imbalance, so that excess buying or selling pressure shows in the
  price;
  \item if still tied, choose the price closest to a reference, such as
  the previous clearing price or an external price feed.
\end{enumerate}

The order of these steps matters. The first two depend on the limit
prices that participants submit, so they are exactly the lever that a
late participant uses in \Cref{prop:lastmover}. Only the third step is
independent of submitted orders. Placing it first would implement the
design response proposed in \Cref{sec:priceimprovement}: it removes the
lever without changing who trades or how much (\Cref{lem:invariance}).

\subsection{Why the Price Is Uniform}

Clearing can also be viewed as an optimisation problem: choose how much
of each order to execute so as to maximise total gains from trade,
subject to the requirement that the quantity bought equals the quantity
sold. In this formulation the uniform price is not an extra rule but
the value of that balance requirement, the price of one additional unit
of the asset in the batch. Every executed order therefore trades at the
same price.

Two consequences follow. First, the optimisation fixes who trades, but
not always the price: several prices can support the same trades, and
this range is precisely $\mathcal{P}^{*}$. The indeterminacy analysed
in \Cref{sec:priceimprovement} is therefore built into the problem, and
the choice within it is exactly what $\sigma$ adds. Second, the fill
logic is simple. Buyers whose limit is above the clearing price and
sellers whose limit is below it are filled in full, and orders on the
wrong side are not filled. Only orders exactly at the clearing price
can be partly filled, and they are the ones subject to rationing.


\subsection{Rationing Rules}
\label{sec:rationing}

At $\clr$ the two sides are in general not equal, $D(\clr) \neq
S(\clr)$, so the orders on the long side with $\pi_k = \clr$ must be
rationed. The rule chosen determines what participants do to improve
their allocation.

\begin{description}
  \item[Price-time priority within the batch.] Orders at $\clr$ are
  filled in submission order. This preserves continuity with continuous
  market conventions and is easy to explain. It also reintroduces a
  speed race at interval scale: if being early inside the batch
  improves expected fill, participants invest in being early, which is
  precisely the incentive the mechanism was built to remove. In the
  terms of \Cref{prop:latency}, it restores a priority payoff to speed
  alongside the deadline payoff that batching already leaves in place.
  \item[Pro rata.] Each order at $\clr$ receives a fill proportional to
  its size; this is the rule adopted, after price priority, in the Dual
  Flow Batch Auction \citep{jump2025dfba}. It is neutral with respect to
  timing, which preserves the core property, but creates a different
  distortion: because allocation increases with submitted size,
  participants submit more than they intend to trade, and the
  equilibrium involves inflated orders that participants expect to be
  scaled down. This is order stuffing, and it degrades the informational
  content of displayed size \citep{haynes2020precedence}.
  \item[Size priority or maximal fill.] Larger orders are filled first,
  or the allocation maximises the number of completely filled orders.
  This favours block liquidity, which may be the intended policy, but
  turns the size-inflation incentive of pro rata into a discrete
  threshold effect.
  \item[Random allocation.] Fills are assigned by lottery among orders
  at $\clr$. This removes both the timing and the size channel, and is
  the only rule under which submitted size is an unbiased statement of
  desired quantity. The cost is execution uncertainty that participants
  cannot manage, which matters for anyone hedging, and the need for a
  randomness source that participants can verify.
  \item[Hybrid rules.] Production systems mix these. The CME's
  K-algorithm family allocates part of the volume by time priority and
  the remainder pro rata, with configurable weights, which lets the
  venue trade off the two distortions rather than choosing one. A
  batch-native analogue would allocate a share $\kappa$ pro rata and
  $1 - \kappa$ randomly, tuning $\kappa$ to balance predictability
  against size inflation.
\end{description}

 No rule is neutral: every rule makes
allocation a function of some variable, and whichever variable that is
becomes the one participants compete on. The design question is
therefore not which rule is undistorted but which distortion the venue
prefers.

% (table unchanged)


\section{Block 5: Settlement and Residual Handling}
\label{sec:settlement}

\subsection{Gross versus Net Settlement}

Once the batch has cleared, the transfers must be executed. Gross
settlement moves each matched pair separately, which is simple and
gives an unambiguous audit trail per fill. Net settlement computes each
participant's net position change and moves only that, which reduces the
number of non-zero transfers, sometimes dramatically when many
participants both buy and sell within the batch. In a decentralised
setting, where each transfer costs gas, the reduction is a direct cost
saving and is one of the main practical arguments for batching
\citep{wang2018loopring,gong2023fbadex}. In a centralised venue the
saving is smaller but the reconciliation burden is still lower. The
trade-off is that netting obscures the correspondence between individual
orders and individual transfers, which complicates per-fill reporting
and any fee schedule levied per execution.

\subsection{The Residual}

The volume on the heavier side that cannot be matched at $\clr$ does not
disappear and is not filled from anywhere else. The design question is
what happens to it, and there are exactly three answers, selected by the
order's time-in-force:

\begin{enumerate}
  \item \textbf{Roll over.} The unmatched remainder is carried into the
  next batch with its original limit price and, where relevant, its
  original submission timestamp. This is the natural default for
  resting limit orders and is the behaviour implemented in this thesis.
  \item \textbf{Cancel.} The remainder is discarded at the end of the
  batch, which is the batch analogue of an immediate-or-cancel order.
  \item \textbf{Cancel the whole order if not fully filled}, the analogue
  of fill-or-kill, which requires the constraint to be imposed inside
  the clearing problem rather than applied afterwards, because an
  all-or-nothing order makes \eqref{eq:lpobj} an integer program.
\end{enumerate}

The third option is worth pausing on, because it shows where the clean
linear structure of the clearing problem breaks. All-or-nothing
constraints turn a problem solvable by a line search over the price grid
into a combinatorial one, which is why venues that support them impose
size limits, and why decentralised batch exchanges that support rich
order types run solver competitions rather than a single deterministic
algorithm \citep{harrison2025hybrid}.

Crucially, no external liquidity source is used to absorb the residual.
The engines compared in \Cref{ch:simulation} can only execute against
order flow that was actually submitted, which is what makes the
comparison like-for-like: neither engine has access to liquidity the
other lacks.
